Положим
\begin{enumerate}
    \item $G^+ := \{z \in \Cm: \im z > 0\}$
    \item $C_R^+ := \{z \in \Cm: |z| = R, \im z \geq 0\}$
    \item $\Cm^{++} := \{z \in \Cm: \im z > 0\}$
    \item $\Cm^{--} := \{z \in \Cm: \im z < 0\}$
\end{enumerate}

\begin{lemma}[<<$+$>>]
	Пусть выполнены следующие условия:
	\begin{enumerate}
		\item Функция $f$ регулярна на $\overline{G}^+ \bs \{z_1, \dotsc, z_n, \infty\}$, где $z_1, \dotsc, z_n \in G^+$~--- ИОТОХ
		\item $\int_{C_R^+}fdz \to 0,\; R \to +\infty$
	\end{enumerate}
	
	Тогда:
	\[\exists\; \text{v.p.} \int_{-\infty}^{\infty} f(x) dx = 2\pi i \sum_{k = 1}^n \res_{z_k}f(z).\]
\end{lemma}

\begin{proof}
	Выберем $R > \max\limits_{k = 1, \ldots, n}|z_k|$, обозначим $D_R = [-R, R] \cup C_R^+$~--- ОСПГ, то есть:	
	\begin{center}
		\scalebox{1}{
			\begin{tikzpicture}
				\clip (-3.4, -0.6) rectangle (3.4, 3);
				\draw [->] (-3, 0) -- (3, 0) node [above, black] {$\re z$};
				\draw [->] (0, -0.5) -- (0, 2.7) node [right, black] {$\im z$};
				
				% A clipped circle is drawn
				\begin{scope}
					\clip (-2.25, 0) rectangle (2.25, 2.2);
					\fill [opacity=0.07] (0, 0) circle[radius=2.2];
					\draw[
					black, line width = 0.2mm,
					decoration={markings, mark=at position 0.13 with {\arrow{>}}},
					decoration={markings, mark=at position 0.38 with {\arrow{>}}},
					postaction={decorate}
					] (0,0) circle[radius=2.2];
					
					\node[draw, star, star points=8, star point ratio=2.2, inner sep=0.5pt, fill, black, label={right, shift={(-3pt,-3pt)}:\scalebox{0.9}{$z_1$}}] at (1.2, 1) {};
					\node[draw, star, star points=8, star point ratio=2.2, inner sep=0.5pt, fill, black, label={right, shift={(-3pt,-3pt)}:\scalebox{0.9}{$z_n$}}] at (-1.5, 0.5) {};
					\node[draw, star, star points=8, star point ratio=2.2, inner sep=0.5pt, fill, black] at (-0.5, 1.4) {};
					\node[draw, star, star points=8, star point ratio=2.2, inner sep=0.5pt, fill, black] at (0.45, 0.5) {};
					\draw (-2.2,0) -- (2.2,0);
				\end{scope}
				
				\draw[line width = 0.2mm, ->] (-2.2,0) -- (-1,0);
				\draw[line width = 0.2mm, ->] (-1,0) -- (1, 0);
				\draw[line width = 0.2mm] (1, 0) -- (2.2,0);
				
				\node[] at (-1.75, 2.1) {$C_{R}^+$};
				\node[] at (0.7, 1.6) {$D_R$};
				
				\draw [black] (2.19,3pt) -- (2.19,-3pt) node [below, black] {$R$};
				\draw [black] (-2.19,3pt) -- (-2.19,-3pt) node [below, black] {$-R$};
			\end{tikzpicture}
		}
	\end{center}
	
	Применим к $D_R$ теорему о вычетах:
	\[\int_{-R}^Rf(x)dx + \int_{C_R^+}fdz = 2\pi i\sum_{k=1}^n\res_{z_k}f\]
	
	Переходя к пределу $R \to +\infty$, получим требуемое.
\end{proof}

\begin{lemma}[<<$-$>>]
	Положим
    \begin{enumerate}
        \item $G^- := \{z \in G: \im z < 0\}$
        \item $C_R^- := \{z \in \Cm: |z| = R, \im z \leq 0\}$.
    \end{enumerate}
    Если выполнено
	\begin{enumerate}
		\item Функция $f$ регулярна на $\overline{G}^- \bs \{z_1, \dotsc, z_n, \infty\}$, где $z_1, \dotsc, z_n \in G^-$~--- ИОТОХ
		\item $\int_{C_R^-}fdz \to 0,\; R \to +\infty$
	\end{enumerate}
	
	Тогда:
	\[\exists\; \text{v.p.} \int_{-\infty}^{\infty} f(x) dx = -2\pi i \sum_{k = 1}^n \res_{z_k}f(z).\]
\end{lemma}

\begin{theorem}
	Пусть выполнены следующие условия:
	\begin{enumerate}
		\item $f = \frac{P}{Q}$, где $P, Q \in \Cm[z]$ "--- многочлены степеней $m, n \in \N \cup \{0\}$ соответственно, не имеющие общих корней
		\item $n \geq m + 2$
		\item $Q$ не имеет вещественных корней и обращается в ноль в точках $z_1, \dotsc, z_N \in \Cm^{++}$ и $w_1, \dotsc, w_M \in \Cm^{--}$
	\end{enumerate}
	
	Тогда верны следующие равенства:
	\[\int_{-\infty}^{+\infty}fdx = 2\pi i\sum_{k=1}^N \res_{z_k}f = -2\pi i\sum_{k=1}^M \res_{w_k}f\]
\end{theorem}

\begin{proof}
	Поскольку $n \ge m + 2$, то существует предел $\lim\limits_{z \to \infty}z^2f(z) = A \in \Cm$. Выберем $R_0$ такое, что при $|z| > R_0$ выполнено неравенство $|z^2f(z)| \le |A| + 1$. Тогда, по свойству оценки интегралов, для любого $R > R_0$ имеем:
	\[\left|\int_{C_R^+}fdz\right| \le (2\pi R)\frac{|A| + 1}{R^2} = 2\pi\frac{|A| + 1}{R} \xrightarrow[R \to \infty]{} 0\]
	
	Значит, выполнены все условия леммы <<$+$>>, из которой следует первое равенство. Второе доказывается аналогично.
\end{proof}

\begin{proposition}[лемма Жордана]
    Положим $G = \{z \in \C | \im z > 0, |z| > R_0\}$, где $R_0 > 0$.
	Пусть выполнены следующие условия:
	\begin{enumerate}
		\item Функция $g$ непрерывна на $G$
		\item Если $M(R) := \max\limits_{z \in C_R^+}|g(z)|$ для любого $R>0$, то $\lim\limits_{R \to +\infty}M(R) = 0$
		\item $\alpha > 0$
	\end{enumerate}
	
	Тогда верно следующее равенство:
	\[\lim_{R \to +\infty} \int_{C_R^+}ge^{i\alpha z}dz = 0\]
\end{proposition}

\begin{proof}
	Произведем замену $z = Re^{it}$, тогда:
	\begin{multline*}
		\left|\int_{C_R^+}g(z)e^{i\alpha z}dz\right| = \left|\int_0^\pi g(Re^{it})e^{i\alpha Re^{it}}z'(t)dt\right| \le M(R)R \int_0^\pi \left|e^{i\alpha R(\cos{t} + i\sin{t})}\right|dt = \\
		= M(R)R \int_0^\pi e^{-\alpha R\sin{t}}dt = 2M(R)R \int_0^{\frac\pi 2} e^{-\alpha R\sin{t}}dt  \stackrel{\sin t \ge \frac 2 \pi t}{\le} 2M(R)R \int_0^{\frac\pi 2} e^{-\alpha \frac{2R}\pi t}dt \le \\
		\le 2M(R)R \int_0^{+\infty} e^{-\alpha \frac{2R}\pi t}dt = 2M(R)R \cdot \frac\pi{2\alpha R} = \frac{M(R)\pi}{\alpha} \xrightarrow[R \to +\infty]{} 0
	\end{multline*}
	
	Получено требуемое.
\end{proof}

\begin{theorem}
	Пусть выполнены следующие условия:
	\begin{enumerate}
		\item $f = \frac PQ$, где $P, Q \in \Cm[z]$ "--- многочлены степеней $m, n \in \N \cup \{0\}$ соответственно, не имеющие общих корней
		\item $n \ge m + 1$
		\item $Q$ не имеет вещественных корней и обращается в ноль в точках $z_1, \dotsc, z_n \in \Cm^{++}$ и $w_1, \dotsc, w_m \in \Cm^{--}$
		\item $\alpha > 0$, $\beta \in \R$
	\end{enumerate}
	
	Тогда верны следующие равенства:
	\[\int_{-\infty}^{+\infty}fe^{i(\alpha x + \beta)}dx = 2\pi i\sum_{k=1}^n \res_{z_k}\left(fe^{i(\alpha x + \beta)}\right) = -2\pi i\sum_{k=1}^m \res_{w_k}\left(fe^{i(\alpha x + \beta)}\right)\]
\end{theorem}

\begin{proof}
	Применим лемму <<$+$>> к функции $fe^{i(\alpha x + \beta)}$. Стремление к нулю интеграла по полуокружности $C_R^+$ достигается засчёт леммы Жордана.
\end{proof}

\begin{unnecessary}
    \begin{note}
    	Если в условиях теоремы выше выполнено $P, Q \in \R[z]$, то для интеграла $I := \int_{-\infty}^{+\infty}fe^{i(\alpha x + \beta)}dx$ выполнены соотношения:
    	\[I = \re I + i \im I = \int_{-\infty}^{+\infty}f\cos{(\alpha x + \beta)}dx + i\int_{-\infty}^{+\infty}f\sin{(\alpha x + \beta)}dx\]
    \end{note}
    
    \begin{note}
    	Пусть $R(u, v)$ "--- рациональная функция. Рассмотрим следующую функцию:
    	\[R_1(z) := R\left(\frac12\left(z + \frac1z\right), \frac1{2i}\left(z - \frac1z\right)\right)\]
    	
    	$R_1$ является рациональной функцией от $z$. Пусть $R_1$ не имеет особых точек на окружности $\{z \in \Cm : |z| = 1\}$. Рассмотрим следующий интеграл:
    	\[I := \int_{|z| = 1}\frac{1}{iz}R_1dz\]
    	
    	С одной стороны, его можно вычислить, используя теорему о вычетах. С другой стороны, если произвести замену $z = e^{it}$, то:
    	\[I = \int_{0}^{2\pi} R(\cos{t}, \sin{t})dt\]
    	
    	Получен способ вычислять еще один класс вещественных интегралов.
    \end{note}
\end{unnecessary}